\section{Matrix models with local trace control}
\label{sec:model}
\label{model:section}

Throughout this section, $A$ is a nonzero separable unital nuclear
$C^*$-algebra; no absorption hypothesis is imposed.  We construct a
matrix model with trace estimates localized to finitely many ideals.  The point is that a norm error
inside a suitable hereditary corner becomes a trace error relative
to the rank of the positive element generating the relevant ideal.
For $k\geq 1$, write
\[
 L_k=\id_A\otimes\mathrm{tr}_k\colon A\otimes M_k\longrightarrow A,
 \qquad
 L_k((z_{rs})_{r,s=1}^k)=\frac1k\sum_{r=1}^k z_{rr},
\]
where $\mathrm{tr}_k$ is the normalized matrix trace.  We may choose
all matrix sizes to be powers of $2$; this convenient restriction
costs nothing in the finite probability-rounding argument below.

Each test uses a positive contraction $a$ and a number $0<\gamma<1$.
We keep the uncut element $a$ for the later comparison target, estimate
the model error on the deeper cut $x=(a-\gamma)_+$, and use the
shallower cut $b=(a-\gamma/2)_+$ to specify both the ideal and the
rank scale for that estimate.  The two cuts have distinct roles:
$x$ lies in the hereditary subalgebra generated by $b$, whose rank
$d_\tau(b)$ will control the trace error.

Choose a countable dense subset $\mathcal D$ of the positive unit ball
of $A$ which is closed under cuts $a\mapsto(a-r)_+$ for positive
rational $r$.  Such a set is obtained from any countable dense subset
by adjoining its successive rational cuts.  Enumerate all pairs
$(a,\gamma)\in\mathcal D\times((0,1)\cap\mathbb Q)$ as
$(a_j,\gamma_j)_{j\geq1}$, and set
\begin{equation}
 \label{model:tests}
 x_j=(a_j-\gamma_j)_+,
 \qquad b_j=(a_j-\gamma_j/2)_+,
 \qquad I_j=\Ideal(b_j).
\end{equation}
In particular, $x_j\in\her(b_j)$ and $d_\tau(x_j)\leq d_\tau(b_j)$
for every lower semicontinuous extended trace $\tau$.

\begin{lemma}[Matrix model]
 \label{model:construction}
There are finite-dimensional $C^*$-algebras $F_n$, powers $k_n$ of
$2$, and completely positive contractive maps
\[
 A\xrightarrow{\alpha_n}F_n
 \xrightarrow{\beta_n}A\otimes M_{k_n},
 \qquad \theta_n=\beta_n\alpha_n,
\]
with the following properties.
\begin{enumerate}
 \item Each $\beta_n$ has order zero, and
 \[
  \|\alpha_n(c)\alpha_n(d)\|\longrightarrow0
  \quad(c,d\in A_+,\ cd=0).
 \]
 \item For every $a\in A$,
 \[
  \|L_{k_n}(\theta_n(a))-a\|\longrightarrow0.
 \]
 \item For every $n,j\geq1$ there is a positive contraction
 $w_{n,j}\in A\otimes M_{k_n}$, and there are nonnegative numbers
 $\eta_{n,j},\delta_{n,j}$ tending to zero as $n\to\infty$ for each
 fixed $j$, such that
 \begin{align}
  \|w_{n,j}-\theta_n(x_j)\|&\leq\eta_{n,j},
       \label{model:retained-norm}\\
  \tau(L_{k_n}(w_{n,j}))
     &\leq \tau(x_j)+\delta_{n,j}d_\tau(b_j)
       \label{model:retained-trace}
 \end{align}
 for every lower semicontinuous extended trace $\tau$ on $A$ with
 $d_\tau(b_j)<\infty$.
\end{enumerate}
The coefficients in \eqref{model:retained-trace} are independent of
$\tau$.  One may take $\delta_{n,j}=2^{-n}$ for all $n,j$.
\end{lemma}

We first establish the finite-stage construction.  Its trace estimates
will be obtained within hereditary subalgebras on which the relevant
extended traces are bounded.

\begin{lemma}[Trace bounds in a hereditary subalgebra]
 \label{model:corner-traces}
Let $P,b\in A_+$ be contractions, and suppose that
$[P]\leq K[b]$ for a finite positive integer $K$.  If $\tau$ is a
lower semicontinuous extended trace with $d_\tau(b)<\infty$, then
its restriction to $D=\her(P)$ extends to a bounded positive linear
functional satisfying
\begin{equation}
 \label{model:corner-bound}
 \|\tau|_D\|\leq d_\tau(P)\leq Kd_\tau(b).
\end{equation}
Consequently, for $y,z\in D_+$,
\begin{equation}
 \label{model:corner-error}
 |\tau(y)-\tau(z)|
    \leq \|y-z\|K d_\tau(b).
\end{equation}
\end{lemma}

\begin{proof}
Choose an approximate unit $(e_l)$ for $D$ consisting of positive
contractive continuous functions of $P$ which vanish at zero.  Such
elements satisfy $d_\tau(e_l)\leq d_\tau(P)$.  For any positive
contraction $z\in D$, the elements $e_lze_l$ converge to $z$ in norm,
and
\[
 \tau(e_lze_l)\leq\tau(e_l^2)\leq d_\tau(P).
\]
Lower semicontinuity therefore gives $\tau(z)\leq d_\tau(P)$.
Monotonicity and additivity of $d_\tau$ on Cuntz classes give the
second inequality in \eqref{model:corner-bound}.  Thus $\tau$ is
bounded on the positive unit ball of $D$ and extends to the asserted
bounded positive functional.  Its norm bound proves
\eqref{model:corner-error}.  This argument also applies when
$d_\tau(b)=0$, in which case the restriction is identically zero.
\end{proof}

\begin{lemma}[Finite ideal slicing]
 \label{model:slicing}
Fix $m,N\geq1$, and let
$\mathcal T=\{1,\ldots,N\}^m$.  For every $t\in\mathcal T$ and
$S\subseteq\{1,\ldots,m\}$, there are positive contractions
$(P_{S,t,i})_{i\geq1}$ in $A$ with the following properties.
\begin{enumerate}
 \item For fixed $t,i$, the family $(P_{S,t,i})_S$ is orthogonal;
 in particular, $\sum_S P_{S,t,i}\leq1$.
 \item For every $a\in A$,
 $\|[P_{S,t,i},a]\|\to0$.  If $1\leq j\leq m$ and $j\notin S$, then
 \[
  \|P_{S,t,i}^{1/2}x_jP_{S,t,i}^{1/2}\|\longrightarrow0.
 \]
 \item If $j\in S$, then $P_{S,t,i}\in I_j$, and there is a finite
 positive integer $K_{S,t,i,j}$ such that
 \begin{equation}
  \label{model:finite-rank-bound}
  [P_{S,t,i}]\leq K_{S,t,i,j}[b_j].
 \end{equation}
 \item For every $a\in A$,
 \begin{equation}
  \label{model:slicing-approximation}
  \limsup_{i\to\infty}
  \left\|\frac1{N^m}\sum_{t\in\mathcal T}\sum_S
     P_{S,t,i}^{1/2}aP_{S,t,i}^{1/2}-a\right\|
  \leq\frac mN\|a\|.
 \end{equation}
\end{enumerate}
\end{lemma}

\begin{proof}
Write $A_\infty=\ell^\infty(A)/c_0(A)$, identifying $A$ with its
constant sequences.  There are commuting positive contractions
$u_1,\ldots,u_m\in A_\infty\cap A'$ such that $u_j$ is represented
by positive contractions in $I_j$ and $u_jx_j=x_j$.  Indeed, at
coordinate $i$ choose successively an element of a quasicentral
approximate unit of each $I_j$.  Require its commutators with the
first $i$ elements of a fixed dense subset of $A$ to have norm less
than $1/i$, require that it act on $x_j$ within $1/i$, and require
that it commute within $1/i$ with the elements already chosen at
that coordinate.  These are finitely many approximate-unit and
quasicentrality requirements.  They give exactly the stated
relations on passing to $A_\infty$.

Choose $N$ pairwise disjoint closed intervals contained in $(0,1)$.
For the $r$th interval choose continuous functions
$g_r^-,g_r^+\colon[0,1]\to[0,1]$ satisfying
\[
 g_r^-g_r^+=0,\qquad
 0\leq g_r^-+g_r^+\leq1,
\]
whose sum is one outside that interval.  We choose $g_r^-$ equal
to one below the interval and zero above it, and $g_r^+$ equal to
zero below the interval and one above it.  Within the interval the
two functions can decrease and increase, respectively, meeting at
zero.  Set $h_r=g_r^-+g_r^+$.  At every point of $[0,1]$ at most
one interval can contribute a loss, so
\begin{equation}
 \label{model:threshold-average}
 1-\frac1N\leq\frac1N\sum_{r=1}^N h_r\leq1.
\end{equation}
For $t=(t_1,\ldots,t_m)\in\mathcal T$ define
\[
 P_{S,t}=
  \prod_{j\in S}g_{t_j}^+(u_j)
  \prod_{j\in\{1,\ldots,m\}\setminus S}g_{t_j}^-(u_j).
\]
These are central positive contractions, and for fixed $t$ they
are pairwise orthogonal as $S$ varies.  Families for different values
of $t$ may overlap and are combined by averaging.  Joint functional
calculus and independent
averaging of the coordinates of $t$ give
\begin{align}
 Q&=\frac1{N^m}\sum_{t\in\mathcal T}\sum_S P_{S,t}
   =\prod_{j=1}^m\left(\frac1N\sum_{r=1}^N h_r(u_j)\right),
       \label{model:mean-slice}\\
 (1-1/N)^m1&\leq Q\leq1,
 \qquad \|1-Q\|\leq m/N.
       \label{model:mean-slice-bound}
\end{align}
If $j\in S$, then $P_{S,t}$ belongs to the image of
$\ell^\infty(I_j)$ in $A_\infty$, since $g_{t_j}^+(0)=0$.
If $1\leq j\leq m$ and $j\notin S$, then
$g_{t_j}^-(u_j)x_j=g_{t_j}^-(1)x_j=0$, whence $P_{S,t}x_j=0$.

We construct the required orthogonal representatives directly,
before imposing exact ideal membership. Start with positive
contractive representatives of a finite orthogonal family, whose
pairwise products tend to zero.  To separate a first representative
$a_i$ from a later representative $b_i$, choose $r_i>0$ tending to
zero slowly enough that $\|a_ib_i\|/r_i\to0$.  Let
$f_i\colon[0,1]\to[0,1]$ be one on $[0,r_i/2]$ and zero on
$[r_i,1]$, with $1-f_i(t)\leq2t/r_i$.  Replace the first element
by $(a_i-r_i)_+$ and every later one by $f_i(a_i)b_if_i(a_i)$.
These replacements are orthogonal to the first cut, and
\[
 \|b_i-f_i(a_i)b_if_i(a_i)\|
    \leq4\|a_ib_i\|/r_i\longrightarrow0.
\]
Choose $r_i$ using the maximum of the finitely many relevant
products.  Repeating with the later elements preserves
orthogonality already obtained and changes no quotient element.
This gives, for every $t$, orthogonal positive contractive lifts
$R_{S,t,i}$ of $P_{S,t}$.

For $j\in S$, the distance from $R_{S,t,i}$ to $I_j$ tends to zero.
Let $\varepsilon_i>0$ tend to zero and dominate all these distances,
over the finite collection of triples $(S,t,j)$ under
consideration.  The first cuts
\[
 T_{S,t,i}=(R_{S,t,i}-2\varepsilon_i)_+
\]
belong exactly to every $I_j$ with $j\in S$: their images in $A/I_j$
vanish by functional calculus.  Choose another sequence
$\sigma_i>0$ tending to zero, and put
\[
 P_{S,t,i}=(T_{S,t,i}-\sigma_i)_+.
\]
Because $T_{S,t,i}\in\Ideal(b_j)$, the usual ideal criterion for
Cuntz comparison gives a finite integer $K_{S,t,i,j}$ for which
\eqref{model:finite-rank-bound} holds.  Here the strict positivity
of the second cut is essential: every positive cut of an element
in $\Ideal(b_j)$ is Cuntz below a finite direct sum of copies of
$b_j$.  Both cuts preserve orthogonality and change the
representatives by at most $2\varepsilon_i+\sigma_i\to0$.
Consequently they preserve all the quotient relations established
above, while imposing all the ideal and finite-multiple conditions
simultaneously.

Finally, centrality of $P_{S,t}$ gives
$P_{S,t}^{1/2}aP_{S,t}^{1/2}=P_{S,t}a$ in $A_\infty$.
Thus \eqref{model:mean-slice-bound} implies
\eqref{model:slicing-approximation}, and $P_{S,t}x_j=0$ proves the
required vanishing of the omitted compressions.  No bound on
$K_{S,t,i,j}$ uniform in $i$ has been asserted or used.
\end{proof}

\begin{lemma}[Finite-stage model]
 \label{model:finite-stage}
Fix $m\geq1$, a finite subset $\mathcal F$ of the unit ball of $A$,
a finite family $\mathcal O$ of pairs of orthogonal positive
contractions in $A$, and $\varepsilon>0$.  There are a
finite-dimensional $C^*$-algebra $F$, a power $k$ of $2$, completely
positive contractive maps
\[
 A\xrightarrow{\alpha}F\xrightarrow{\beta}A\otimes M_k,
 \qquad\theta=\beta\alpha,
\]
and positive contractions $w_1,\ldots,w_m\in A\otimes M_k$ such
that $\beta$ has order zero and
\begin{align}
 \|L_k(\theta(a))-a\|&<\varepsilon
       &&(a\in\mathcal F), \label{model:finite-approx}\\
 \|\alpha(c)\alpha(d)\|&<\varepsilon
       &&((c,d)\in\mathcal O), \label{model:finite-orth}\\
 \|w_j-\theta(x_j)\|&<\varepsilon
       &&(1\leq j\leq m), \label{model:finite-retained}\\
 \tau(L_k(w_j))&\leq\tau(x_j)+\varepsilon d_\tau(b_j)
       &&(1\leq j\leq m, d_\tau(b_j)<\infty).
       \label{model:finite-trace}
\end{align}
The last inequality holds for every lower semicontinuous extended
trace $\tau$ satisfying the indicated finite-rank condition.
\end{lemma}

\begin{proof}
\emph{Fixing the corners.}
Choose $N$ so large that $m/N<\varepsilon/4$, and use
Lemma~\ref{model:slicing}.  Write
\[
 l_{S,t,i}(a)=P_{S,t,i}^{1/2}aP_{S,t,i}^{1/2}.
\]
Asymptotic centrality implies, for $(c,d)\in\mathcal O$,
\[
 \|l_{S,t,i}(c)l_{S,t,i}(d)\|
 =\|P_{S,t,i}^{1/2}cP_{S,t,i}dP_{S,t,i}^{1/2}\|
 \longrightarrow0.
\]
Together with \eqref{model:slicing-approximation} and the omitted
compression estimates, this permits us to fix one coordinate $i$
such that
\begin{align}
 \left\|\frac1{N^m}\sum_{t,S}l_{S,t,i}(a)-a\right\|
   &<\varepsilon/3 &&(a\in\mathcal F),
       \label{model:fixed-average}\\
 \|l_{S,t,i}(x_j)\|&<\varepsilon
       &&(1\leq j\leq m,\ j\notin S), \label{model:fixed-omitted}\\
 \|l_{S,t,i}(c)l_{S,t,i}(d)\|&<\varepsilon/2
       &&((c,d)\in\mathcal O). \label{model:fixed-products}
\end{align}
All these requirements range over finite sets.  From now on this
coordinate is fixed.  For each occurrence indexed by $(S,t)$,
suppress $i$ and write
\[
 P=P_{S,t,i},\quad D_P=\her(P),\quad
 l_P(a)=P^{1/2}aP^{1/2},\quad K_{P,j}=K_{S,t,i,j}.
\]
Every subscript $P$ below retains this $(S,t)$ occurrence index.
Thus coincident positive elements, including zero pieces, still
receive separate corner approximations and outcome labels.
In particular, the finitely many constants $K_{P,j}$ have already
been fixed before any approximation accuracy within $D_P$ is
chosen.

\emph{Approximating within the corners.}
Each $D_P$ is separable and nuclear.  The Brown--Carri\'on--White
approximation theorem, Theorem~\ref{prelim:bcw}, supplies
finite-dimensional approximations with asymptotically order-zero
first maps and return maps which are finite convex combinations
of completely positive contractive order-zero maps.
For a fixed corner, take a sequence of these first maps
$\zeta_r\colon D_P\to E_r$. The asymptotic product estimate
\eqref{prelim:asymptotic-product} gives
\begin{equation}
 \label{model:quotient-product-bound}
 \limsup_{r\to\infty}\|\zeta_r(y)\zeta_r(z)\|
   \leq\|yz\|\qquad(y,z\in D_P).
\end{equation}
Thus sufficiently late first maps have product norm less than
$\varepsilon$ on the finitely many pairs
$(l_P(c),l_P(d))$ in \eqref{model:fixed-products}.  This uses the
approximation sequence for the fixed corner; it requires no
uniformity in a varying hereditary subalgebra.

Choose positive numbers $\xi_P$ such that
\begin{equation}
 \label{model:inner-tolerances}
 \max_P\xi_P<\varepsilon/3,
 \qquad
 \frac1{N^m}\sum_{t\in\mathcal T}\sum_{S\ni j}
       K_{P_{S,t,i},j}\xi_{P_{S,t,i}}<\varepsilon/2
       \quad(1\leq j\leq m).
\end{equation}
There are only finitely many constants in these requirements, so
such choices are possible.  Choose a corresponding approximation
in each $D_P$, denoted
\[
 D_P\xrightarrow{\zeta_P}E_P
   \xrightarrow{\varphi_P}D_P,
 \qquad
 \varphi_P=\sum_{r=1}^{h_P}\lambda_{P,r}\beta_{P,r},
\]
where the $\lambda_{P,r}$ are nonnegative and sum to one, and each
$\beta_{P,r}$ is completely positive contractive and order zero.
All component return maps for this corner have the same domain
$E_P$ and use the same first map $\zeta_P$.  Arrange that
\begin{align}
 \|\varphi_P\zeta_Pl_P(a)-l_P(a)\|&<\xi_P
       &&\bigl(a\in\mathcal F\cup\{x_1,\ldots,x_m\}\bigr),
       \label{model:inner-approx}\\
 \|\zeta_Pl_P(c)\zeta_Pl_P(d)\|&<\varepsilon
       &&((c,d)\in\mathcal O).
       \label{model:inner-products}
\end{align}
For a zero corner take $E_P=0$ and the zero maps, with a single
return component of weight one.

\emph{Individual outcomes and their retained parts.}
Consider the finite probability space whose outcomes $\nu$
consist of a tuple $t\in\mathcal T$ and, for each subset $S$, a
choice $r_S\in\{1,\ldots,h_{P_{S,t,i}}\}$.  Give $t$ uniform
probability and make the return-map choices independently with
the probabilities $\lambda_{P,r}$.  Write $\lambda_\nu$ for
the resulting probability, and let
\[
 E_t=\bigoplus_S E_{P_{S,t,i}},\qquad
 \alpha_t(a)=\bigoplus_S\zeta_{P_{S,t,i}}l_{P_{S,t,i}}(a),
\]
\[
 \beta_\nu((z_S)_S)=
   \sum_S\beta_{P_{S,t,i},r_S}(z_S),
 \qquad T_\nu=\beta_\nu\alpha_t.
\]
The map $\alpha_t$ is completely positive contractive because a
finite direct sum has the maximum norm.  The hereditary
subalgebras $D_{P_{S,t,i}}$ are mutually orthogonal for the single
tuple $t$ used by the outcome.
Consequently $\beta_\nu$ is completely positive contractive
and order zero: each summand has these properties, and products
between distinct summand ranges vanish.  Moreover,
\eqref{model:inner-products} bounds the product defect of every
$\alpha_t$ by $\varepsilon$, uniformly over the outcomes.

Independence of the return-map choices gives
\begin{equation}
 \label{model:expected-map}
 \sum_\nu\lambda_\nu T_\nu(a)
   =\frac1{N^m}\sum_{t\in\mathcal T}\sum_S
       \varphi_{P_{S,t,i}}\zeta_{P_{S,t,i}}l_{P_{S,t,i}}(a).
\end{equation}
For each fixed $t$, the differences in
\eqref{model:inner-approx} lie in orthogonal hereditary
subalgebras, so their sum has norm at most $\max_P\xi_P$.
Equations \eqref{model:fixed-average},
\eqref{model:inner-tolerances}, and \eqref{model:expected-map}
therefore imply
\begin{equation}
 \label{model:expected-approximation}
 \left\|\sum_\nu\lambda_\nu T_\nu(a)-a\right\|
    <2\varepsilon/3\qquad(a\in\mathcal F).
\end{equation}

For $1\leq j\leq m$, retain only the corners indexed by $S\ni j$:
\begin{equation}
 \label{model:outcome-retained}
 W_{\nu,j}=
    \sum_{S\ni j}\beta_{P_{S,t,i},r_S}
                \zeta_{P_{S,t,i}}l_{P_{S,t,i}}(x_j).
\end{equation}
This is a positive contraction.  Contractivity of each component
and orthogonality of the omitted corners give, for every
individual outcome,
\begin{equation}
 \label{model:before-averaging}
 \|T_\nu(x_j)-W_{\nu,j}\|
   \leq\max_{S\not\ni j}\|l_{P_{S,t,i}}(x_j)\|
   <\varepsilon.
\end{equation}
The maximum over an empty family is taken to be zero.  In
particular, the estimate holds before taking any average.

\emph{The uniform relative trace estimate.}
Fix $j$ and a lower semicontinuous extended trace $\tau$ with
$d_\tau(b_j)<\infty$.  In every retained corner,
Lemma~\ref{model:corner-traces} applies with the already fixed
constant $K_{P,j}$.  Thus \eqref{model:inner-approx} implies
\[
 \tau(\varphi_P\zeta_Pl_P(x_j))
    \leq\tau(l_P(x_j))+\xi_PK_{P,j}d_\tau(b_j)
    \qquad(j\in S).
\]
All quantities in this inequality are finite.  For each fixed
$t$, traciality and $\sum_{S\ni j}P_{S,t,i}\leq1$ give
\begin{align*}
 \sum_{S\ni j}\tau(l_{P_{S,t,i}}(x_j))
 &=\tau\left(x_j^{1/2}
       \Bigl(\sum_{S\ni j}P_{S,t,i}\Bigr)x_j^{1/2}\right)\\
 &\leq\tau(x_j).
\end{align*}
Taking expectations in \eqref{model:outcome-retained} and using
\eqref{model:inner-tolerances} proves
\begin{equation}
 \label{model:expected-trace}
 \sum_\nu\lambda_\nu\tau(W_{\nu,j})
    \leq\tau(x_j)+(\varepsilon/2)d_\tau(b_j).
\end{equation}
For later use there is also an individual bound
\begin{equation}
 \label{model:individual-trace}
 \tau(W_{\nu,j})\leq M_{\nu,j}d_\tau(b_j),
 \qquad M_{\nu,j}=\sum_{S\ni j}K_{P_{S,t,i},j}<\infty.
\end{equation}
Every coefficient here was chosen independently of $\tau$.
If $d_\tau(b_j)=0$, then each retained output has trace exactly
zero, as does $x_j$.  Hence the same inequalities hold in this
case without an absolute additive error.  No condition is being
placed on traces for which $d_\tau(b_j)=\infty$.

\emph{Rounding the probabilities and forming the matrix model.}
After all corner maps and the constants $M_{\nu,j}$ have been
fixed, choose $k=2^s$ and nonnegative integers $n_\nu$ with
$\sum_\nu n_\nu=k$ so that
$\mu_\nu=n_\nu/k$ satisfies
\begin{equation}
 \label{model:rounding-errors}
 \sum_\nu|\mu_\nu-\lambda_\nu|<\varepsilon/3,
 \qquad
 \sum_\nu|\mu_\nu-\lambda_\nu|M_{\nu,j}
      <\varepsilon/2\quad(1\leq j\leq m).
\end{equation}
Such choices exist for arbitrarily large $s$: round the finitely
many numbers $k\lambda_\nu$ down to integers and distribute
the remaining units among the outcomes.  The resulting
probabilities stay in the probability simplex, and their total
absolute error tends to zero as $s\to\infty$.  Contractivity of
the $T_\nu$ shows that replacing $\lambda_\nu$ by
$\mu_\nu$ changes the expected map on the unit ball by less
than $\varepsilon/3$.  Equations
\eqref{model:individual-trace} and \eqref{model:rounding-errors}
show that the expected retained trace changes in absolute value
by at most $(\varepsilon/2)d_\tau(b_j)$.

List the outcomes with their multiplicities as
$\nu_1,\ldots,\nu_k$, and write $t_l$ for the threshold
tuple in $\nu_l$.  Define
\[
 F=\bigoplus_{l=1}^k E_{t_l},\qquad
 \alpha(a)=\bigoplus_{l=1}^k\alpha_{t_l}(a),\qquad
 \beta((z_l)_{l=1}^k)
    =\operatorname{diag}\bigl(\beta_{\nu_l}(z_l)\bigr)_{l=1}^k.
\]
The map $\alpha$ is completely positive contractive, and $\beta$
is completely positive contractive and order zero, by the
orthogonality of the matrix diagonal positions.  We have
\[
 \theta(a)=\operatorname{diag}(T_{\nu_l}(a))_{l=1}^k,
 \qquad
 L_k(\theta(a))=\sum_\nu\mu_\nu T_\nu(a).
\]
Consequently \eqref{model:expected-approximation} and
\eqref{model:rounding-errors} give \eqref{model:finite-approx},
and the maximum norm in the direct sum gives
\eqref{model:finite-orth}.  Finally, put
\[
 w_j=\operatorname{diag}(W_{\nu_l,j})_{l=1}^k.
\]
These are positive contractions, and the estimate for individual
outcomes in \eqref{model:before-averaging} proves
\eqref{model:finite-retained}.  Since
$L_k(w_j)=\sum_\nu\mu_\nu W_{\nu,j}$,
\eqref{model:expected-trace} and the trace part of the rounding
estimate prove \eqref{model:finite-trace}.
\end{proof}

\begin{proof}[Proof of Lemma~\ref{model:construction}]
Fix a countable dense subset $(c_l)_{l\geq1}$ of the unit ball of
$A$.  The space
\[
 \{(c,d)\in A_+\times A_+:
           \|c\|,\|d\|\leq1,\ cd=0\}
\]
is a separable metric space, so choose a countable dense family
$(c'_l,d'_l)_{l\geq1}$ in this space.  At stage $n$, apply
Lemma~\ref{model:finite-stage} to the first $n$ tests,
$\mathcal F_n=\{c_1,\ldots,c_n\}$,
$\mathcal O_n=\{(c'_1,d'_1),\ldots,(c'_n,d'_n)\}$, and
$\varepsilon=2^{-n}$.  Denote the resulting maps by $\alpha_n$
and $\beta_n$, and the retained elements by $w_{n,j}$ for
$j\leq n$.

Contractivity extends \eqref{model:finite-approx} from the dense
family to point-norm convergence on $A$.  Likewise, if $c,d$ are
orthogonal positive contractions and $(c',d')$ is one of the
chosen orthogonal pairs, then
\[
 \|\alpha_n(c)\alpha_n(d)\|
 \leq\|c-c'\|+\|d-d'\|
       +\|\alpha_n(c')\alpha_n(d')\|.
\]
First taking the limit superior in $n$, then approximating
$(c,d)$ by such pairs, proves asymptotic preservation of
orthogonality.  Scaling proves the assertion for arbitrary
orthogonal positive elements.

For $j>n$ set $w_{n,j}=0$.  Take $\delta_{n,j}=2^{-n}$ for all
$n,j$, and set
\[
 \eta_{n,j}=
 \begin{cases}
  2^{-n},&j\leq n,\\
  1,&j>n.
 \end{cases}
\]
Equations \eqref{model:retained-norm} and
\eqref{model:retained-trace} now hold for every $n,j$, and both
error sequences tend to zero for each fixed $j$.  At each stage
the trace inequality holds simultaneously for all the indicated
extended traces.  Thus the diagonalization involves only the
countable lists of elements, tests, and orthogonal pairs; it
requires neither a countable family of traces nor a restriction
on the ideal lattice of $A$.
\end{proof}

\begin{remark}
 \label{model:composite-order-zero}
The composites $\theta_n$ also asymptotically preserve
orthogonality.  Indeed, the order-zero product estimate gives
\[
 \|\theta_n(c)\theta_n(d)\|
 \leq\|\alpha_n(c)\alpha_n(d)\|
 \qquad(c,d\in A_+).
\]
Consequently they induce a c.p.c. order-zero map into
$\prod_n(A\otimes M_{k_n})/\bigoplus_n(A\otimes M_{k_n})$,
and into the corresponding norm ultraproduct.  The same conclusion
holds after composing the $\beta_n$ with any c.p.c. order-zero maps
into a fixed target algebra: the displayed product estimate remains
valid for each composite return map.
\end{remark}
